\documentclass[11pt]{article}
\usepackage[margin=1.15in]{geometry}
\usepackage{amsmath,amssymb,amsthm}
\usepackage{booktabs,array}
\usepackage[colorlinks=true,linkcolor=blue,citecolor=blue,urlcolor=blue,
  pdftitle={Constraint Inventory for fm-exact.py}]{hyperref}

\newtheorem{lemma}{Lemma}[section]
\newtheorem{proposition}[lemma]{Proposition}
\theoremstyle{remark}
\newtheorem{remark}[lemma]{Remark}
\newtheorem{convention}[lemma]{Convention}

\newcommand{\N}{N}
\newcommand{\eof}[1]{e\!\left(#1\right)}
\newcommand{\code}[1]{\texttt{#1}}

\title{Constraint Inventory for the Exact Certification\\
of the Level Profile $\zeta(a)=(5a-4)/4$\\[2mm]
\large Companion documentation to Appendix~A\\
of ``An Explicit Semiprime-Tail Bound in Square-Root Intervals\\
and a One-Hypothesis Reduction for $[4n^2-n,\,4n^2+n]$''}
\author{Michael M. Ross\\
\textit{Independent Researcher}\\
\small \texttt{michaelmross@cantab.net}}

\begin{document}
\maketitle

\begin{abstract}\noindent
This note is the full constraint inventory referred to in Appendix~A of the
paper referenced above (henceforth [P]). It derives every linear inequality appearing
in the certification script \code{fm\_exact.py} from the published theorem
that inequality encodes; states the sign, negation, and boundary conventions
needed to read the code's coefficient tuples; records the three places
where the encoding is a deliberate \emph{relaxation} of the true tool, and
proves that the certification remains sound there by a sandwich against the
Coverage Lemma of [P]; and cross-references every atom family to the code
identifiers and to the corresponding statements of [P]. This is designed to be an audit trail, not new mathematics.
\end{abstract}

\section{Purpose and provenance}\label{sec:purpose}

The script \code{fm\_exact.py} certifies, in exact rational arithmetic,
the level profile
\[
\zeta(\gamma)\;=\;\frac{6-5\gamma}{4}\qquad(\gamma\in[1,6/5]),
\]
in both directions: for $\beta<\zeta(\gamma)$ every estimation block
admits an admissible configuration (achievability --- proved by hand in
[P], Coverage Lemma, and re-verified here as a byproduct), and at
$\beta=\zeta(\gamma)$ a block exists that no encoded tool covers
(optimality --- [P], Optimality Proposition). The certification is a
finite computation because every tool's admissibility region is a
conjunction of inequalities that are \emph{linear} in the block
coordinates at fixed $\gamma$. This document is the dictionary between
those inequalities, the code, and the literature.

\section{Conventions}\label{sec:conventions}

\begin{convention}[coordinates and exponents]\label{conv:coords}
A slice is fixed by $\gamma\in[1,6/5]$ (prime block $p\sim
\N^{\gamma}$; equivalently $a=2-\gamma$, cofactor $q\asymp\N^{a}$). The
unknowns are $x=(\mu,\eta,\beta)$: the Vaughan-piece exponent, the dyadic
frequency exponent, and the sieve-level exponent. All exponents are with
respect to base $\N=2n$; $x_{\mathrm{arith}}:=4n^2\asymp\N^2$. We write
$\Sigma:=\beta+\gamma$ and $\xi:=\eta+2-\Sigma$ (phase-size exponent).
The admissibility requirement for a block, [P] eq.\ (trilinear
requirement), is that the estimating bound have exponent at most
$\Sigma$ (i.e.\ relative saving $\N^{\eta}$ against the trivial size
$\N^{\eta+\Sigma}$), with an $x^{-c\varepsilon}$ margin suppressed
throughout this document: all atoms are stated with $\varepsilon=0$ and
\emph{closed} inequalities. Facet equalities are therefore boundary
points of both the pass and fail regions; this is intentional (see
Convention~\ref{conv:boundary}).
\end{convention}

\begin{convention}[atom format in the code]\label{conv:atoms}
In \code{failure\_atoms()} a constraint is the tuple
$\code{C}(a_{\mu},a_{\eta},a_{\beta},c_0,c_{\gamma})$, meaning
\[
a_{\mu}\,\mu+a_{\eta}\,\eta+a_{\beta}\,\beta\ \le\ c_0+c_{\gamma}\,\gamma ,
\]
all entries rational. Each \emph{tool} passes on a conjunction of pass
conditions; the code stores, per tool, the list of \emph{negated} pass
conditions (``fail atoms''), one of which must hold for the tool to fail
at a point. A \emph{failure witness} is a point where every tool fails;
distributing the disjunctions gives the product over tools of the
fail-atom counts:
\[
\underbrace{3}_{\code{FI0}}\times
\underbrace{3}_{\code{FI1}}\times
\underbrace{3}_{\code{FI2}}\times
\underbrace{2}_{\code{RS0}}\times
\underbrace{2}_{\code{RS1}}\times
\underbrace{2}_{\code{RS2}}\times
\underbrace{1}_{\code{RS3}}\times
\underbrace{1}_{\code{ELEM}}\;=\;216
\]
candidate polytopes per piece type (\code{ELEM} is present only for
type~I, where its single fail option is the conjunction of its three
negated clauses; type~II has $216$ combinations without the final
factor). Each polytope comprises the chosen fail atoms plus the region
constraints of Section~\ref{sec:region} ($\approx14$ inequalities in $3$
unknowns); the minimal $\beta$ over it is found by exhaustive vertex
enumeration ($\binom{14}{3}=364$ candidate bases, $3\times3$ systems
solved by Cramer's rule over $\mathbb{Q}$ in \code{solve3()}), and
$\zeta(\gamma)$ is the minimum over all polytopes (\code{min\_beta()}).
\end{convention}

\begin{convention}[boundaries and monotonicity]\label{conv:boundary}
Because all fail atoms are closed, the computed $\zeta(\gamma)$ is the
infimum of levels admitting a witness, and the witness printed at
$\beta=\zeta(\gamma)$ sits on the facet with equalities active.
Failure is not a priori monotone in $\beta$ (raising $\beta$ widens the
frequency range and tightens windows, but also widens the reachable
grouping intervals), so the script's output is, by itself, only the
statement ``the \emph{first} failure occurs at $\beta=\zeta(\gamma)$.''
That everything strictly below the facet passes is the Coverage Lemma of
[P], proved by hand; the two artifacts meet at the facet. See also
Section~\ref{sec:soundness}.
\end{convention}

\section{The block model and the region $R$}\label{sec:region}

Per [P] (Decomposition and grouping), a block carries variables
$(h,e_1,e_2,m_v,k)$ with exponents
$(\eta,\ \beta-s,\ s,\ \mu,\ \gamma-\mu)$, $s\in[0,\beta]$ free. Vaughan
pieces: type~I has $\mu\in[0,2\gamma/3]$ with $k$ smooth; type~II has
$\mu\in[\gamma/3,\gamma/2]$, both coefficients arbitrary. The frequency
range from the pipeline lemma is $0\le\eta\le\Sigma-1$ (blocks with
$\Sigma\le1$ require nothing). The region constraints in
\code{region()}:
\begin{center}
\begin{tabular}{@{}llll@{}}
\toprule
code tuple & inequality & role & piece type\\
\midrule
\code{C(-1,0,0,0,0)} & $\mu\ge0$ & piece range & I\\
\code{C(1,0,0,0,2/3)} & $\mu\le 2\gamma/3$ & piece range & I\\
\code{C(-1,0,0,0,-1/3)} & $\mu\ge\gamma/3$ & piece range & II\\
\code{C(1,0,0,0,1/2)} & $\mu\le\gamma/2$ & piece range & II\\
\code{C(0,-1,0,0,0)} & $\eta\ge0$ & frequency range & both\\
\code{C(0,1,-1,-1,1)} & $\eta\le\beta+\gamma-1$ & frequency range & both\\
\code{C(0,0,-1,0,0)}, \code{C(0,0,1,1,0)} & $0\le\beta\le1$ & search box & both\\
\bottomrule
\end{tabular}
\end{center}

\section{Tool family RS: the Robert--Sargos theorem}\label{sec:RS}

\subsection{Statement and specialization}
Robert--Sargos \cite[Theorem~1]{RS}: for exponents with
$\alpha(\alpha-1)\beta_0\gamma_0\ne0$, dyadic ranges $H,N_1,M_1$, phase
size $X\ge1$, and arbitrary coefficients $|a(h,\nu)|\le1$, $|b(w)|\le1$,
\[
\sum_{h\sim H}\sum_{\nu\sim N_1}a(h,\nu)\sum_{w\sim M_1}b(w)\,
\eof{X\tfrac{h^{\beta_0}\nu^{\gamma_0}w^{\alpha}}
{H^{\beta_0}N_1^{\gamma_0}M_1^{\alpha}}}
\ \ll_{\varepsilon}\ (HN_1M_1)^{1+\varepsilon}
\Bigl\{\Bigl(\tfrac{X}{HN_1M_1^{2}}\Bigr)^{1/4}
+(HN_1)^{-1/4}+M_1^{-1/2}+X^{-1/2}\Bigr\}.
\]
We use $(\beta_0,\gamma_0,\alpha)=(1,-1,-1)$, so
$\alpha(\alpha-1)\beta_0\gamma_0=-2\ne0$; the standalone slot $w$ has
exponent $m_1$, the joint slot $(h,\nu)$ has exponents
$(\eta,\Sigma-m_1)$, and $X=\N^{\xi}$ with $\xi\ge2-\Sigma\ge3/4>0$.

\subsection{Window derivation}\label{ssec:RSwindow}
The requirement is $(\eta+\Sigma)+\max_i t_i\le\Sigma$, i.e.\
$\max_i t_i\le-\eta$, where in exponents, with $n_1=\Sigma-m_1$ and
$\xi-\eta=2-\Sigma$:
\begin{align*}
t_1&=\tfrac{\xi-\eta-n_1-2m_1}{4}=\tfrac{2-2\Sigma-m_1}{4}
&&\Rightarrow\quad m_1\ \ge\ 4\eta+2-2\Sigma
&&\text{(A)}\\
t_2&=-\tfrac{\eta+n_1}{4}
&&\Rightarrow\quad m_1\ \le\ \Sigma-3\eta
&&\text{(B)}\\
t_3&=-\tfrac{m_1}{2}
&&\Rightarrow\quad m_1\ \ge\ 2\eta
&&\text{(C)}\\
t_4&=-\tfrac{\xi}{2}
&&\Rightarrow\quad \eta\ \le\ 2-\Sigma
&&\text{(D)}
\end{align*}
Subsumptions used by the encoding: (A)$\,\Leftarrow\,$(C) since
$4\eta+2-2\Sigma\le2\eta\iff\eta\le\Sigma-1$, which is a region
constraint; (D) is implied on the relevant part of the region (see
Remark~\ref{rem:relaxations}). The effective window is
$m_1\in[2\eta,\ \Sigma-3\eta]$, matching [P] (Admissibility window
Lemma), whose window is nonempty iff $5\eta\le\Sigma$.

\subsection{The four grouping families and their code atoms}
The reachable standalone exponents are
$m_1\in\{s,\ \mu+s,\ \gamma-\mu+s,\ \gamma+s\}$, $s\in[0,\beta]$
([P], available-set display; the fourth family is retained here because
the RS window's upper cap $\Sigma-3\eta$ can exceed $\gamma$ at small
$\eta$). Taking in each family the choice $m_1=\max(\text{base},2\eta)$
where possible, the pass conditions and their negations:

\begin{center}
\small
\begin{tabular}{@{}lllll@{}}
\toprule
tool & pass conditions & fail atoms (code) & fail inequality\\
\midrule
\code{RS0} ($m_1=s$) & $2\eta\le\beta$;\ $5\eta\le\Sigma$
 & \code{C(0,-2,1,0,0)} & $\beta\le2\eta$\\
 & & \code{C(0,-5,1,0,-1)} & $\beta\le5\eta-\gamma$\ \ (i.e.\ $5\eta\ge\Sigma$)\\
\code{RS1} ($m_1=\mu+s$) & $\mu+\beta\ge2\eta$;\ $\mu\le\Sigma-3\eta$
 & \code{C(1,-2,1,0,0)} & $\mu+\beta\le2\eta$\\
 & & \code{C(-1,-3,1,0,-1)} & $\mu\ge\Sigma-3\eta$\\
\code{RS2} ($m_1=\gamma-\mu+s$) & $\gamma-\mu+\beta\ge2\eta$;\ $\gamma-\mu\le\Sigma-3\eta$
 & \code{C(-1,-2,1,0,-1)} & $\gamma-\mu+\beta\le2\eta$\\
 & & \code{C(1,-3,1,0,0)} & $\mu+\beta\le3\eta$\ \ (i.e.\ $\gamma-\mu\ge\Sigma-3\eta$)\\
\code{RS3} ($m_1=\gamma+s$) & $3\eta\le\beta$\ \ (i.e.\ $\gamma\le\Sigma-3\eta$)
 & \code{C(0,-3,1,0,0)} & $\beta\le3\eta$\\
\bottomrule
\end{tabular}
\end{center}

\begin{remark}[documented relaxations in the RS encoding]
\label{rem:relaxations}
Three omissions make the encoded tools slightly \emph{stronger} than the
true ones on part of the search box:
(i) \code{RS1} and \code{RS2} omit the window-nonemptiness condition
$5\eta\le\Sigma$ (which \code{RS0} includes);
(ii) \code{RS3} omits the lower-edge condition $\gamma+s\ge2\eta$, vacuous
only when $2\eta\le\gamma$, i.e.\ when $\Sigma\le(\gamma+2)/2$;
(iii) condition (D) of Subsection~\ref{ssec:RSwindow} is omitted
everywhere, valid only when $\eta\le2-\Sigma$, i.e.\ when
$\Sigma\le3/2$ given $\eta\le\Sigma-1$.
All three omitted conditions \emph{hold automatically} on
$\{\Sigma\le5/4,\ \eta\le\Sigma-1\}$, which contains the entire facet
region; outside it they are genuine relaxations. Soundness is unaffected:
see Section~\ref{sec:soundness}.
\end{remark}

\section{Tool family ELEM: elementary bounds (type~I only)}
\label{sec:elem}

For a type~I block the smooth variable $k\sim K=\N^{\kappa}$,
$\kappa=\gamma-\mu$, is summed with the phase $f(k)=hz'/(ck)$,
$c=e_1e_2m_v\sim\N^{\chi}$, $\chi=\beta+\mu$, all other variables in
absolute value; $|f'|\asymp\Lambda=\N^{\lambda}$,
$\lambda=\eta+2-\chi-2\kappa$. Three clauses, whose disjunction is the
pass condition; the tool fails only if all three fail, which is the
single \code{ELEM} conjunction in the code.

\begin{itemize}
\item[(E1)] \emph{Kuzmin--Landau branch} ($\lambda<0$): by
Kuzmin--Landau (e.g.\ \cite[Theorem~2.1]{GK}, after dyadic subdivision
so that $|f'|\le\tfrac12$), $\sum_k\eof{f(k)}\ll\Lambda^{-1}$; total
exponent $\eta+\chi-\lambda=2(\chi+\kappa)-2=2\Sigma-2\le\Sigma$ for
$\Sigma\le2$: \emph{always admissible in region}. Pass clause:
$\lambda<0\iff\eta\le\beta-\mu+2\gamma-2$.\\
Fail atom \code{C(-1,-1,1,2,-2)}: $\ \eta\ge\beta-\mu+2\gamma-2$.
\item[(E2)] \emph{Second-derivative test / exponent pair
$(\tfrac12,\tfrac12)$} \cite[Theorem~2.2]{GK}:
$\sum_k\ll(\Lambda K)^{1/2}+(K/\Lambda)^{1/2}$. The first term needs
$\eta+\chi+\tfrac{\lambda+\kappa}{2}\le\Sigma$, which simplifies to
$3\eta+2\mu\le\beta+3\gamma-2$; the companion term needs
$\eta+\chi+\tfrac{\kappa-\lambda}{2}=\tfrac{\eta+3\Sigma-2}{2}\le\Sigma$,
i.e.\ $\eta\le2-\Sigma$, automatic on the facet region.\\
Fail atom \code{C(-2,-3,1,2,-3)}: $\ 3\eta+2\mu\ge\beta+3\gamma-2$.
\item[(E3)] \emph{Exponent pair $(\tfrac16,\tfrac23)$}
\cite[\S3]{GK}: $\sum_k\ll\Lambda^{1/6}K^{2/3}(+\,\text{admissible
companions as in (E1),(E2)})$; the main term needs
$\eta+\chi+\tfrac{\lambda}{6}+\tfrac{2\kappa}{3}
=\tfrac{7\eta+5\beta+3\mu+2\gamma+2}{6}\le\Sigma$, which simplifies to
$7\eta+3\mu\le\beta+4\gamma-2$.\\
Fail atom \code{C(-3,-7,1,2,-4)}: $\ 7\eta+3\mu\ge\beta+4\gamma-2$.
\end{itemize}

\begin{remark}
[P] uses only (E1) and (E2) in the Coverage Lemma; (E3) is carried in
the certification solely so that the Optimality Proposition quantifies
over the larger toolkit. The certified optimum is identical with or
without (E3) --- one of the two byproducts recorded in [P],
Appendix~A.
\end{remark}

\section{Tool family FI: Fouvry--Iwaniec Theorem~6}\label{sec:FI}

This family is \emph{not used} in the proofs of [P]; it is encoded so
that the Optimality Proposition's claim --- the profile is unchanged
when \cite[Theorem~6]{FI} is adjoined --- is machine-checked rather
than asserted. Since the derivation appears nowhere in [P], it is
recorded here in full.

\subsection{Statement and exponent form}
\cite[Theorem~6]{FI}: for $\alpha\ne0,1$, an additive character
$\chi(h)$, and arbitrary $|\varphi_m|,|\psi_n|\le1$,
\begin{align*}
S_{\chi\varphi\psi}(H,M,N)\ \ll\ (HMN)^{1/2}\Bigl[\,&(H+N)^{1/2}
\bigl(x_6^{1/8}H^{-1/6}M^{1/12}N^{1/6}
+x_6^{1/8}H^{-1/8}N^{3/8}\\
&\qquad+N^{1/2}+N^{1/4}M^{1/8}\bigr)\,x_6^{1/8}
\;+\;M^{1/2}\;+\;x_6^{-1/4}M^{1/2}N\,\Bigr]\mathcal{L}^4,
\end{align*}
where $x_6$ is the phase size. With slot exponents $(m_0,n_0)$,
$m_0+n_0=\Sigma$, frequency $\eta$, and $x_6=\N^{\xi}$, the requirement
(bound $\le\N^{\Sigma}$) yields, after the prefactor
$(\eta+\Sigma)/2$:
\begin{itemize}
\item the $M^{1/2}$ term (T5): $\tfrac{\eta+\Sigma}{2}+\tfrac{m_0}{2}
\le\Sigma\iff\eta\le n_0$ --- a \emph{lower} window edge, which also
fixes $(H+N)^{1/2}\mapsto\N^{n_0/2}$ in the remaining terms;
\item five upper caps on $n_0$, derived by direct substitution
($\xi=\eta+2-\Sigma$, $m_0=\Sigma-n_0$):
\[
\begin{array}{lll}
\text{T1 } (x_6^{1/4}H^{-1/6}M^{1/12}N^{1/6}\cdot):
& 7n_0\le 8\Sigma-7\eta-6
& \Rightarrow\ n_0\le\mathrm{cap}_{1}=\tfrac{8\Sigma-7\eta-6}{7}\\[1mm]
\text{T2 } (x_6^{1/4}H^{-1/8}N^{3/8}\cdot):
& 7n_0\le 6\Sigma-5\eta-4
& \Rightarrow\ n_0\le\mathrm{cap}_{2}=\tfrac{6\Sigma-5\eta-4}{7}\\[1mm]
\text{T3 } (x_6^{1/8}N^{1/2}\cdot):
& 8n_0\le 5\Sigma-5\eta-2
& \Rightarrow\ n_0\le\mathrm{cap}_{3}=\tfrac{5\Sigma-5\eta-2}{8}\\[1mm]
\text{T4 } (x_6^{1/8}N^{1/4}M^{1/8}\cdot):
& 5n_0\le 4\Sigma-5\eta-2
& \Rightarrow\ n_0\le\mathrm{cap}_{4}=\tfrac{4\Sigma-5\eta-2}{5}\\[1mm]
\text{T6 } (x_6^{-1/4}M^{1/2}N):
& 2n_0\le 2-\Sigma-\eta
& \Rightarrow\ n_0\le\mathrm{cap}_{6}=\tfrac{2-\Sigma-\eta}{2}
\end{array}
\]
\end{itemize}
(Sample derivation, T2: multiply the requirement
$\tfrac{\eta+\Sigma}{2}+\tfrac{n_0}{2}+\tfrac{\xi}{4}-\tfrac{\eta}{8}
+\tfrac{3n_0}{8}\le\Sigma$ by $8$ and substitute $\xi$:\ 
$5\eta+7n_0+2\Sigma+4\le8\Sigma$. The others are identical in kind.)

\subsection{Pruning lemma}\label{ssec:pruning}
\begin{lemma}
On $\{\eta\le\Sigma-1,\ \Sigma\le4/3\}$:
$\mathrm{cap}_2\le\mathrm{cap}_1$,
$\mathrm{cap}_2\le\mathrm{cap}_3$, and
$\mathrm{cap}_2\le\mathrm{cap}_4$. Hence the FI window is
$[\,\eta,\ \min(\mathrm{cap}_6,\mathrm{cap}_2)\,]$.
\end{lemma}
\begin{proof}
$\mathrm{cap}_1-\mathrm{cap}_2=\tfrac{2\Sigma-2\eta-2}{7}\ge0$ iff
$\eta\le\Sigma-1$.
$\mathrm{cap}_2\le\mathrm{cap}_4$ iff $5(6\Sigma-5\eta-4)\le
7(4\Sigma-5\eta-2)$ iff $\Sigma+5\eta\le3$; with $\eta\le\Sigma-1$ this
follows from $6\Sigma\le8$, i.e.\ $\Sigma\le4/3$.
$\mathrm{cap}_2\le\mathrm{cap}_3$ iff $8(6\Sigma-5\eta-4)\le
7(5\Sigma-5\eta-2)$ iff $13\Sigma\le5\eta+18$, implied by
$\Sigma\le4/3<18/13$.
\end{proof}

\subsection{Grouping families and code atoms}
The reachable $n_0$ families are $\{s,\ \mu+s,\ \gamma-\mu+s\}$,
$s\in[0,\beta]$; the fourth family $n_0=\gamma+s$ is \emph{provably
dead} for this tool, since $\gamma\ge1$ while
$\mathrm{cap}_6\le\tfrac{2-\Sigma}{2}\le\tfrac12$, and is therefore not
encoded. In each family we take $n_0$ at the base value (lower edge
condition base${}+\beta\ge\eta$):

\begin{center}
\small
\begin{tabular}{@{}llll@{}}
\toprule
tool & pass conditions & fail atoms (code) & fail inequality\\
\midrule
\code{FI0} & $\eta\le\beta$;\ $\eta\le\mathrm{cap}_6$;\ $\eta\le\mathrm{cap}_2$
 & \code{C(0,-1,1,0,0)} & $\eta\ge\beta$\\
 & & \code{C(0,-3,-1,-2,1)} & $3\eta+\beta\ge2-\gamma$\ \ ($\eta\ge\mathrm{cap}_6$)\\
 & & \code{C(0,-12,6,4,-6)} & $12\eta-6\beta\ge6\gamma-4$\ \ ($\eta\ge\mathrm{cap}_2$)\\
\code{FI1} & $\mu+\beta\ge\eta$;\ $\mu\le\mathrm{cap}_6$;\ $\mu\le\mathrm{cap}_2$
 & \code{C(1,-1,1,0,0)} & $\mu+\beta\le\eta$\\
 & & \code{C(-2,-1,-1,-2,1)} & $2\mu+\eta+\beta\ge2-\gamma$\\
 & & \code{C(-7,-5,6,4,-6)} & $7\mu+5\eta-6\beta\ge6\gamma-4$\\
\code{FI2} & $\gamma-\mu+\beta\ge\eta$;\ $\gamma-\mu\le\mathrm{cap}_6$;\
$\gamma-\mu\le\mathrm{cap}_2$
 & \code{C(-1,-1,1,0,-1)} & $\gamma-\mu+\beta\le\eta$\\
 & & \code{C(2,-1,-1,-2,3)} & $2\mu-\eta-\beta\le3\gamma-2$\\
 & & \code{C(7,-5,6,4,1)} & $7\mu-5\eta+6\beta\le\gamma+4$\\
\bottomrule
\end{tabular}
\end{center}

\noindent(The last two \code{FI2} rows are the sign-flipped forms of
$\gamma-\mu\ge\mathrm{cap}_6$ and $\gamma-\mu\ge\mathrm{cap}_2$.) The
pruning is a further relaxation outside $\Sigma\le4/3$, handled by the
same soundness argument below. The trivial-size lower edge $\eta\le
n_0$ is encoded as the first pass condition in each family.

\section{Soundness of the relaxed encoding}\label{sec:soundness}

Let $\zeta_{\mathrm{code}}(\gamma)$ be the script's output and
$\zeta_{\mathrm{true}}(\gamma)$ the first-failure level of the
\emph{exactly} encoded toolkit. Every deviation documented in
Remark~\ref{rem:relaxations} and Subsection~\ref{ssec:pruning} makes an
encoded tool pass on a superset of its true pass region; failure
witnesses are therefore harder to find, so
$\zeta_{\mathrm{code}}\ge\zeta_{\mathrm{true}}$. In the other
direction, the Coverage Lemma of [P] proves by hand --- using only the
exactly-valid parts of the encoding, all active inside
$\{\Sigma\le5/4\}$ where every omitted condition holds --- that all
blocks pass for $\beta<(6-5\gamma)/4$, so
$\zeta_{\mathrm{true}}\ge(6-5\gamma)/4$. Since the script returns
$\zeta_{\mathrm{code}}=(6-5\gamma)/4$ exactly at every tested rational
$\gamma$, the sandwich closes:
\[
\frac{6-5\gamma}{4}\;=\;\zeta_{\mathrm{code}}
\;\ge\;\zeta_{\mathrm{true}}\;\ge\;\frac{6-5\gamma}{4}.
\]
Finally, the printed witnesses are legitimate witnesses for the
\emph{true} system: at $(\mu,\eta)=(\gamma/2,\Sigma-1)$ with
$4\beta+5\gamma=6$ one checks directly that every omitted condition
($5\eta\le\Sigma$, $\eta\le2-\Sigma$, $2\eta\le\gamma$,
$\Sigma\le4/3$) is satisfied, so no relaxed tool was wrongly credited
there.

\section{Cross-reference table}\label{sec:xref}

\begin{center}
\small
\begin{center}
\small
\begin{tabular}{@{}ll>{\raggedright\arraybackslash}p{0.52\textwidth}@{}}
\toprule
code identifier & this document & item in [P]\\
\midrule
\code{RS0}--\code{RS3} & \S\ref{sec:RS} & \S5.4 Admissibility-window Lemma (5.3);
\S5.3 available-set display (10)\\
\code{ELEM} clause (E1) & \S\ref{sec:elem} & \S5.5 elementary Lemma (5.4),
Kuzmin--Landau term $\Lambda^{-1}$\\
\code{ELEM} clause (E2) & \S\ref{sec:elem} & \S5.5 elementary Lemma (5.4),
exponent pair $(\tfrac12,\tfrac12)$; used in \S5.6 Coverage Lemma (5.5)\\
\code{ELEM} clause (E3) & \S\ref{sec:elem} & \S5.6 Optimality Prop.\ (5.6);
Appendix~A (certification only)\\
\code{FI0}--\code{FI2} & \S\ref{sec:FI} & \S5.6 Optimality Prop.\ (5.6);
Appendix~A byproducts (certification only)\\
\code{region()} & \S\ref{sec:region} & \S5.3 Decomposition and grouping;
\S5.1 pipeline Lemma (5.1), frequency range below (9)\\
\code{solve3()}, \code{min\_beta()} & Convention~\ref{conv:atoms} &
Appendix~A (vertex enumeration)\\
\bottomrule
\end{tabular}
\end{center}
\end{center}

\section{Reproduction}\label{sec:repro}

Run \code{python3 fm\_exact.py} (Python~3; \code{numpy}/\code{scipy}
are used only for the final floating-point root-finding of
$\theta^{\ast\ast}$, which is verification of a closed-form calculus step and
not part of the certification). Expected output: the exact table
$\zeta(\gamma)=1/4,\,3/16,\,1/8,\,1/16,\,0$ at
$\gamma=1,\,21/20,\,11/10,\,23/20,\,6/5$, agreement with
$(6-5\gamma)/4$ throughout, the witness
$(\mu,\eta)=(11/20,\,9/40)$ at $\gamma=11/10$, the threshold
$\theta^{\ast\ast}=0.939305$, and the margins $+0.0136$, $+0.2028$,
$+0.5478$ (in units of $u$) at $\vartheta=0.94,\,0.95,\,0.97$.
Runtime is a few minutes.

\subsection*{Acknowledgments}
The author thanks Claude Fable 5.1 (Anthropic) for collaboration in analysis and exposition.

\begin{thebibliography}{9}
\bibitem{FI} E.~Fouvry and H.~Iwaniec, \emph{Exponential sums with
monomials}, J.~Number Theory \textbf{33} (1989), 311--333.
\bibitem{GK} S.~W. Graham and G.~Kolesnik, \emph{Van der Corput's
Method of Exponential Sums}, London Math.\ Soc.\ Lecture Note Ser.\
\textbf{126}, Cambridge Univ.\ Press, 1991.
\bibitem{RS} O.~Robert and P.~Sargos, \emph{Three-dimensional
exponential sums with monomials}, J.~reine angew.\ Math.\ \textbf{591}
(2006), 1--20.
\end{thebibliography}

\end{document}